\originalsection{作业 2: 牛顿定律}
\begin{exercise}
原题 1. 上块 $m_1=2.2$ kg 放在下块 $m_2=4.8$ kg 上, 桌面光滑, 块间静摩擦系数 $\mu_s=0.5$. 水平力作用在上块, 求两块不相对滑动的最大力 (N).
\end{exercise}
\begin{solution}
共同加速度 $a=F/(m_1+m_2)$. 下块仅受水平摩擦, 所需 $f=m_2a\le\mu_sm_1g$. 故 $F_{\max}=\mu_sm_1g(m_1+m_2)/m_2=15.74$ N. 力若施于下块, 阈值不同, 因而不能忽略作用对象.
\end{solution}
\begin{exercise}
原题 2. 卡车在水平直线匀加速 $a$, 车后挂无质量不可伸长绳和质量 $M$ 的桶. 初始晃动后桶与车相对位置固定. 绳与车后门 (竖直方向) 夹角 $\theta$. (a) 求 $\theta$. (b) 求张力. (c)(d) 卡车改在向下倾角 $\alpha$ 的坡上沿坡继续加速 $a$, 后门垂直坡面, 求新的夹角及张力.
\end{exercise}
\begin{solution}
(a)(b) 惯性系中 $T\sin\theta=Ma$, $T\cos\theta=Mg$, 所以 $\tan\theta=a/g$, $T=M\sqrt{g^2+a^2}$. (c)(d) 沿下坡及法向投影, 有效重力切向分量为 $g\sin\alpha-a$, 向坡内法向分量 $g\cos\alpha$. 绳对桶必须反向, 故相对后门有向夹角满足 $\tan\theta=(a-g\sin\alpha)/(g\cos\alpha)$ (正向约定为桶落在车后). 张力 $T=M\sqrt{(a-g\sin\alpha)^2+g^2\cos^2\alpha}$. 若 $a<g\sin\alpha$, 桶偏向车前, 无符号角应取绝对值并另说明方向.
\end{solution}
\begin{exercise}
原题 3. 倾角 $\theta$ 的楔紧靠墙静止, 小箱质量 $m$ 沿楔滑下. 箱、楔、地面、墙面均无摩擦. (a) 箱所受合力大小是多少? (b) 墙对楔的力大小是多少?
\end{exercise}
\begin{solution}
法向力 $N=mg\cos\theta$, 箱的切向合力为 $mg\sin\theta$, 所以 (a) $F_{\rm total}=mg\sin\theta$. (b) 箱对楔的水平分量为 $N\sin\theta$, 墙平衡该量, $F_{ww}=mg\sin\theta\cos\theta$.
\end{solution}
\begin{exercise}
原题 4. 上块 $m_2$ 叠在下块 $m_1$ 上, 两块右端由轻绳绕右方固定无摩擦轮连接. 所有接触面动摩擦系数均为 $\mu$. 水平向左拉上块 $m_2$ 使它匀速左移, 求拉力 $F$.
\end{exercise}
\begin{solution}
绳长不变使下块匀速向右. 下块受向左的块间摩擦 $\mu m_2g$ 及地面摩擦 $\mu(m_1+m_2)g$, 张力向右, 故 $T=\mu g(m_1+2m_2)$. 上块受张力和向右摩擦 $\mu m_2g$, 因此 $F=\mu g(m_1+3m_2)$. 两块间的相对速度决定块间摩擦方向.
\end{solution}
\begin{exercise}
原题 5. (a) 弹簧 $k$ 挂质量 $m$, 伸长 $d$. 同样弹簧下分两根绳分别挂 $m/2$, 求伸长 $x$. (b) 高度 $h$ 的平台质量 $M$, 中心有支柱, 两侧对称连接两根向地面弹簧 $k_1$ 与两根向天花板弹簧 $k_2$. 有柱时均无形变. 去柱后求平衡高度 $h_f$.
\end{exercise}
\begin{solution}
(a) 弹簧承受总重量相同, $x=d=mg/k$. (b) 平台下移 $\delta$ 使下方弹簧压缩、上方弹簧拉长, 四根恢复力合为 $2(k_1+k_2)\delta$. 平衡给 $\delta=Mg/[2(k_1+k_2)]$, $h_f=h-\delta$. 需 $\delta<h$ 且弹簧可承受相应拉压, 否则平台碰地或线性模型失效.
\end{solution}
